% Generated draft: requires mathematical and source review.
\documentclass[11pt]{article}
\usepackage[margin=25mm]{geometry}
\usepackage{fontspec}
\setmainfont{Latin Modern Roman}
\usepackage{amsmath,amssymb,mathtools,unicode-math}
\setmathfont{Latin Modern Math}
\usepackage{enumitem,needspace,xurl,hyperref}
\setlength{\parindent}{0pt}
\setlength{\parskip}{5pt}
\setlength{\emergencystretch}{3em}
\newcommand{\N}{\mathbb{N}}
\newcommand{\Z}{\mathbb{Z}}
\newcommand{\Q}{\mathbb{Q}}
\newcommand{\R}{\mathbb{R}}
\newcommand{\C}{\mathbb{C}}
\newcommand{\F}{\mathbb{F}}
\newcommand{\A}{\mathbb{A}}
\newcommand{\PP}{\mathbb{P}}
\newcommand{\E}{\mathbb{E}}
\newcommand{\eps}{\varepsilon}
\newcommand{\dd}{\,\mathrm{d}}
\begin{document}
\section*{UnsolvedMath ID: 30003899}
\textbf{Source problem code:} catalogue entry
\subsection{Definitions and mathematical statement}
Let \(g\geq 1\) be an integer. A principally polarized abelian variety \((A,\Theta)\) of dimension \(g\) consists of an abelian variety \(A\) of dimension \(g\) and an effective ample Cartier divisor \(\Theta\) such that, for \(L=\mathcal O_A(\Theta)\), the polarization morphism
\[
\phi_L:A\longrightarrow A^\vee,\qquad a\longmapsto t_a^*L\otimes L^{-1},
\]
is an isomorphism. Here \(A^\vee\) is the dual abelian variety and \(t_a\) is translation by \(a\). The divisor \(\Theta\) is called a theta divisor. The principally polarized variety is indecomposable if it is not isomorphic, as a polarized abelian variety, to a product of two principally polarized abelian varieties of positive dimension.

For a point \(x\in\Theta\), let \(f\in\mathcal O_{A,x}\) be a local equation of \(\Theta\), and let \(\mathfrak m_x\) be the maximal ideal of the regular local ring \(\mathcal O_{A,x}\). Define
\[
\operatorname{mult}_x\Theta
=\max\{m\in\Z_{\geq 0}:f\in\mathfrak m_x^m\}.
\]
This order is independent of the choice of local equation.

\textbf{Proposed bound.} For every integer \(g\geq1\), every indecomposable principally polarized abelian variety \((A,\Theta)\) of dimension \(g\), and every point \(x\in\Theta\), is it true that
\[
\operatorname{mult}_x\Theta\leq\frac{g+1}{2}?
\]
The quantifier includes points on non-isolated singular loci; no isolated-singularity hypothesis appears in the supplied statement.

\textbf{Unresolved scope.} The source record does not specify a ground field, its characteristic, whether indecomposability is geometric, or whether points are closed or geometric points. The displayed assertion is consequently a statement schema pending these conventions, rather than a fully specified assertion over an invented ground field.
\subsection{Short English statement}
The proposed bound asks whether every point of a theta divisor on an indecomposable principally polarized abelian variety of dimension \(g\) has multiplicity at most \((g+1)/2\), including points on non-isolated singular loci. The ground field and point conventions remain unspecified. The supplied record labels the problem partially solved; this draft does not independently verify that assessment.
\subsection{Sources}
\begin{itemize}
\item Local registry: \texttt{data/unsolvedmath/problems.json}, with \texttt{catalog/problems.json} as the supplementary registry; supplied permanent UnsolvedMath ID \(30003899\), title \emph{Multiplicity Bounds for Principal Theta Divisors}, exact source code \texttt{catalogue entry}. The record belongs to the Oberwolfach Reports Open Problems collection. No individual proposer is identified; the supplied proposed year is 2018.
\item Supplied report citation: \emph{Classical Algebraic Geometry (2019)}, Oberwolfach Reports, DOI 10.4171/owr/2018/33, pp. 1983--2031. \url{https://doi.org/10.4171/owr/2018/33}. The background describes Oberwolfach Reports as documenting mathematical workshops and their open-problem discussions.
\item Supplied literature reference, labelled Mustata--Popa: \url{https://arxiv.org/abs/1605.08088}.
\item Supplied literature reference, labelled Hodge ideals: \url{https://people.math.harvard.edu/~mpopa/papers/SI.pdf}.
\item Supplied literature reference, labelled OWR report: \url{https://ems.press/content/serial-article-files/46758?nt=1}.
\item Literature assessment dated exactly \texttt{2026-08-21}: the supplied curation labels the problem partially solved and states that the bound is proved for isolated singularities and other special cases, while the conjecture for arbitrary principally polarized abelian varieties with non-isolated singularities remains open. This is supplied dated curation derived from the OWR record, not an independently verified literature assessment. Its phrase ``arbitrary principally polarized abelian varieties'' omits the indecomposability hypothesis present in the problem statement and does not justify broadening that statement. The record expressly calls for independent verification of time-sensitive status and cited claims.
\item The supplied statement status is \texttt{corrected\_verified}; its accompanying explanation says that mathematical content was preserved while notation, typography, and omitted definitions were repaired. This is a supplied editorial designation, not mathematical verification performed for this draft.
\end{itemize}
\end{document}
